\documentclass[conference]{IEEEtran}
\IEEEoverridecommandlockouts

\usepackage{cite}
\usepackage{amsmath,amssymb,amsfonts}
\usepackage{algorithmic}
\usepackage{graphicx}
\usepackage{textcomp}
\usepackage{xcolor}
\usepackage{booktabs}
\usepackage{multirow}
\usepackage{url}
\usepackage{microtype}

\def\BibTeX{{\rm B\kern-.05em{\sc i\kern-.025em b}\kern-.08em
    T\kern-.1667em\lower.7ex\hbox{E}\kern-.125emX}}

\begin{document}

\title{Evaluating the Transferability of Carbon-Aware Cloud Scheduling\\
Across Electricity Grids: A Time-Aware Study of\\
Great Britain and India}

\author{
    \IEEEauthorblockN{Vignesh Jaisankar, Kopathy K}
    \IEEEauthorblockA{
        School of Computer Science and Engineering (SCOPE)\\
        Vellore Institute of Technology, Chennai, India
    }
}

\maketitle

\begin{abstract}
Carbon-aware cloud scheduling reduces operational emissions by aligning batch
workload execution with periods of lower grid carbon intensity. Most published
evaluations assess scheduling policies within a single electricity grid,
leaving open whether reported gains transfer across grids with different
generation mixes and diurnal patterns. This paper presents a controlled,
time-aware comparative evaluation of seven scheduling policies across two
national grids: Great Britain~(GB) and India. Indian source timestamps,
recorded in Indian Standard Time~(IST, UTC+05:30), are explicitly converted to
UTC and resampled onto a common hourly timeline before any cross-region
comparison, removing a potential 5.5-hour alignment error present in earlier
versions of this work. Using a five-fold temporal cross-validation scheme with
four workload seeds per fold, we evaluate the Bahreini SRR-APX scheduler,
three Carbon-Adaptive Deadline-Sensitive~(CADS) configurations, a
deadline-bounded carbon-shifting policy, the Wiesner~2021 Ad-Hoc policy, and a
first-come first-served~(FCFS) baseline. For SRR-APX, mean carbon savings
relative to FCFS were 6.04\% in GB (95\%~CI: 2.66\%--9.41\%) and 6.16\% in
India (95\%~CI: 4.93\%--7.38\%), a difference of 0.12 percentage points, with
India marginally higher. SRR-APX achieves zero material deadline misses in both
regions while admitting approximately 93--95\% of the sixty submitted jobs. The
India policy comparison reveals a pronounced tension between carbon savings and
service quality: the Wiesner~2021 Ad-Hoc policy achieves 5.78\% savings but
incurs 35.30 material misses per run on average, against zero for SRR-APX. A
separate three-seed forecast-guided investigation yields a small descriptive
improvement of 0.42 percentage points, though the limited seed count and
identical outcomes across forecast and oracle variants preclude strong
conclusions. The findings underscore that regional performance must be measured
rather than assumed, and that maximising carbon savings in isolation can
materially degrade service quality.
\end{abstract}

\begin{IEEEkeywords}
carbon-aware scheduling, cloud computing, electricity carbon intensity,
deadline-constrained scheduling, green computing, temporal workload shifting,
IST--UTC alignment
\end{IEEEkeywords}

%--------------------------------------------------------------
\section{Introduction}
\label{sec:intro}
%--------------------------------------------------------------

Data centres account for a substantial share of global electricity
consumption, with recent estimates placing them at approximately 200--250~TWh
per year and growing \cite{masanet2020recalibrating,iea2022datacenters}. The
carbon intensity of grid electricity, measured in grams of
CO\textsubscript{2}~equivalent per kilowatt-hour~(gCO\textsubscript{2}e/kWh),
varies significantly across hours of the day and across national grids,
depending on the prevailing generation mix and dispatch order
\cite{radovanovic2023,siler2012should}. Carbon-aware scheduling exploits this
variability by deferring flexible workloads to periods when the grid is
comparatively clean, reducing emissions without altering computation outcomes
\cite{radovanovic2023,acun2023}.

Prior work has demonstrated that temporal workload shifting can yield
measurable carbon savings in specific regions
\cite{wiesner2021,hanafy2024gaia,bahreini2023dispatcher}. A recurring
methodological limitation, however, is that evaluations are typically conducted
within a single grid or a narrow set of North American or Western European
contexts \cite{frachtenberg2023limitations}. Electricity grids differ
substantially in their generation mix, diurnal patterns, and seasonal
characteristics. India's grid remains heavily coal-dependent at the national
level, with renewable penetration growing but not yet dominating the daily
carbon curve \cite{cea2024,ceai2023general}. Great Britain's grid has a
proportionally larger share of wind and nuclear, producing more pronounced
intra-day carbon variability \cite{neso2024api}. The assumption that a
scheduler performing well in one region will transfer its gains to another is
therefore not self-evident and deserves direct empirical examination.

A second methodological concern specific to multi-region studies is temporal
alignment. Indian electricity authority data is published in Indian Standard
Time~(IST, UTC+05:30). If source timestamps are ingested without timezone
conversion, workload events recorded in UTC are paired with carbon-intensity
observations that are 5 hours and 30 minutes misaligned. Earlier iterations
of this work identified this risk; the corrected pipeline described here
explicitly interprets India source timestamps as IST, converts them to UTC, and
resamples observations onto a common hourly timeline. Residual provenance
uncertainty is acknowledged in Section~\ref{sec:threats}.

Service quality must be evaluated alongside carbon savings. Carbon-aware
scheduling commonly introduces waiting time: workloads are deferred to cleaner
windows, extending the gap between submission and completion. For jobs with
explicit deadlines, this creates the risk of material violations. Some prior
analyses report savings figures without adequately distinguishing policies that
achieve them at modest service cost from those that impose heavy deadline
penalties \cite{hanafy2024gaia,bahreini2022approx}. A policy achieving 9\%
savings but missing deadlines on one-third of runs is not equivalent, in
deployment terms, to one achieving 6\% savings with zero material misses.

This paper addresses these gaps through two experimental investigations.
Experiment~A is a controlled, time-aware comparative evaluation of seven
scheduling policies across GB and India. Experiment~B is a separate,
exploratory investigation of forecast-guided variants of the Bahreini SRR-APX
scheduler using India carbon-intensity data. The two experiments use distinct
workload configurations and evaluation protocols and are presented separately
throughout.

The paper makes the following contributions:

\begin{enumerate}
  \item A corrected IST-to-UTC timestamp alignment pipeline for Indian CEA
        generation data, removing a 5.5-hour misalignment present in earlier
        versions of this work.
  \item A time-aware, five-fold comparative evaluation of seven scheduling
        policies across GB and India, using 20 paired workload runs per
        region-policy combination (Experiment~A).
  \item A joint analysis of carbon savings, job admission rate, completion
        rate, material deadline misses, waiting time, and scheduler runtime
        across all seven policies in both regions.
  \item Characterisation of the carbon-savings--service-quality trade-off
        frontier, showing that SRR-APX achieves the highest savings in both
        regions among policies with zero material misses.
  \item A separate three-seed forecast-guided investigation~(Experiment~B)
        reporting a small descriptive improvement from seasonal-naive forecast
        weighting, with the oracle-guided diagnostic producing identical
        aggregate outcomes.
\end{enumerate}

The remainder of this paper is organised as follows. Section~\ref{sec:related}
reviews related work. Section~\ref{sec:method} describes the methodology.
Section~\ref{sec:results} presents the results. Section~\ref{sec:discussion}
discusses the findings. Section~\ref{sec:threats} addresses threats to validity.
Section~\ref{sec:conclusion} concludes.

%--------------------------------------------------------------
\section{Related Work}
\label{sec:related}
%--------------------------------------------------------------

\subsection{Carbon-Aware Computing and Temporal Shifting}

The principle of aligning computing workloads with cleaner periods of grid
electricity was operationalised at scale in Google's Carbon-Intelligent
Computing System~\cite{radovanovic2023}, which uses virtual capacity curves and
day-ahead carbon forecasts to shift temporally flexible datacenter tasks.
Wiesner~et~al.~\cite{wiesner2021} provided simulation-based quantification of
temporal shifting across four European and North American regions, analysing
the joint influence of deadline slack, scheduling strategy, and forecast error.
Both works established that available savings are tightly coupled to
carbon-intensity variability in the target grid.

Hanafy~et~al.~\cite{frachtenberg2023limitations} subsequently analysed
temporal and spatial workload shifting across 123 grid regions, finding that
practical savings are limited for many clouds and workloads and that simple
greedy policies frequently capture most of the available gain. Acun~et~al.\
\cite{acun2023} proposed Carbon Explorer as a holistic framework for
carbon-aware datacenter design, modelling both operational and embodied
components. Together these results temper expectations about large-percentage
savings from scheduling alone.

\subsection{Deadline-Aware and Service-Quality Constraints}

Hanafy~et~al.~\cite{hanafy2024gaia} introduced GAIA, addressing a three-way
trade-off among carbon emissions, job completion time, and cloud cost for batch
workloads. Hanafy~et~al.~\cite{hanafy2024carbonscaler} proposed CarbonScaler,
exploiting workload elasticity to vary server allocation in response to carbon
signals while maintaining throughput. Both works measure deadline or
completion-time overheads alongside carbon savings, a practice this paper
follows. Lechowicz~et~al.~\cite{lechowicz2023pause} studied the online pause
and resume problem and applied it to carbon-aware load shifting under
deadline constraints.

LP-relaxation-based approaches with randomised rounding have a formal basis
in the scheduling and energy-efficiency literature. Bampis~et~al.\
\cite{bampis2014energy} proposed a unifying configuration-LP framework with
randomised rounding for energy-minimisation problems on heterogeneous
processors. The Bahreini SRR-APX scheduler extends this framework with
carbon-intensity-dependent objective coefficients~\cite{bahreini2023dispatcher,
bahreini2022approx}.

\subsection{Geographical and Regional Grid Variation}

Comparative evaluations across national grids are less common than
within-region studies. Wiesner~et~al.~\cite{wiesner2021} compared four
regions but focused on temporal shifting under forecast error rather than
cross-region transferability of a specific scheduler. Souza~et~al.\
\cite{souza2023casper} introduced CASPER for scheduling across geographically
distributed cloud regions with different carbon intensities. Hanafy~et~al.\
\cite{frachtenberg2023limitations} explicitly modelled the role of regional
carbon variability in limiting achievable savings. Neither of those studies,
however, evaluates the service-quality behaviour of a specific scheduler across
a coal-heavy Asian grid and a more variable European grid under a controlled
temporal validation scheme, nor does either address the IST-to-UTC alignment
problem that arises when using published Indian electricity authority data.

\begin{figure}[t]
    \centering
    \includegraphics[width=\columnwidth]{paper_diurnal_local_clock.pdf}
    \caption{Mean diurnal carbon-intensity profiles by source-local clock:
Great Britain in UTC and India in IST. Both profiles are derived from
the carbon-intensity traces used in the scheduling evaluation.}
    \label{fig:carbon-profiles}
\end{figure}

\subsection{Forecast Methods for Carbon Intensity}

Forecast accuracy affects temporal shifting performance. Radovanovic~et~al.\
\cite{radovanovic2023} used day-ahead carbon-intensity forecasts in Google's
operational system. Seasonal-naive forecasting, which repeats the observation
from the corresponding period in the prior cycle, is a standard low-complexity
baseline for short-to-medium-range time-series prediction~\cite{makridakis2020m4,
hyndman2021forecasting}. It is used in Experiment~B as the deployable forecasting
method, in contrast to the oracle-guided diagnostic, which uses future
carbon-intensity data and is not operationally deployable.

Cross-validation methodology for time-series data requires care to avoid
overly optimistic estimates from non-sequential splits. Bergmeir~et~al.\
\cite{bergmeir2018note} establish conditions under which sequential blocked
folds provide valid error estimates; this paper uses five sequential temporal
folds in accordance with that principle.

%--------------------------------------------------------------
\section{Methodology and Experimental Design}
\label{sec:method}
%--------------------------------------------------------------

\subsection{Research Questions and Experimental Structure}

This paper addresses two research questions. The primary question
(Experiment~A) asks how scheduling effectiveness changes when the same seven
policies are applied to carbon-intensity patterns from Great Britain and India,
evaluated jointly across carbon savings, job admission, material deadline
misses, waiting time, and scheduler runtime. The secondary question
(Experiment~B) asks whether adding seasonal-naive forecast guidance to the
Bahreini SRR-APX randomised-rounding procedure changes scheduling outcomes.
The two experiments use distinct workload generators and evaluation protocols;
their results are never merged or ranked jointly.

\begin{figure*}[t]
    \centering    \includegraphics[width=0.6\textwidth]
    {ecocloud_architecture.pdf}
    \caption{Overall architecture of the EcoCloud carbon-aware
    scheduling evaluation. Regional carbon-intensity traces are
    preprocessed and evaluated across five temporal folds using
    multiple scheduling policies. A separate forecast-guided
    experiment compares forecast-based scheduling with an
    Oracle diagnostic that uses future carbon-intensity data.}
    \label{fig:architecture}
\end{figure*}

\subsection{Carbon-Intensity Data Sources and Preprocessing}
\label{subsec:data}

\subsubsection{Great Britain}

Carbon-intensity data for GB was obtained from the NESO Carbon Intensity
API~\cite{neso2024api}, which provides half-hourly actual and forecast
carbon-intensity values in gCO\textsubscript{2}e/kWh at national granularity.
GB timestamps include explicit UTC offsets, supporting unambiguous temporal
alignment with the experiment start. Half-hourly observations are aggregated
to hourly means to match the common timeline.

\subsubsection{India}
\label{subsubsec:india-data}

Carbon-intensity data for India was derived from generation-mix data published
by the Central Electricity Authority~(CEA)~\cite{cea2024,ceai2023general} and
processed to hourly carbon-intensity estimates. The preprocessing pipeline
performs the following steps:

\begin{enumerate}
  \item Source CSV timestamps are parsed as IST-localised datetime objects
        (e.g., \texttt{2026-08-19~00:00:00~Asia/Kolkata}).
  \item Each timestamp is converted to UTC by subtracting the IST offset of
        5~hours and 30~minutes, yielding UTC-normalised observations.
  \item The UTC series is resampled to a uniform hourly grid; missing
        single-hour gaps are filled by forward-fill and longer gaps are
        filled by linear interpolation.
  \item The resulting UTC-aligned series is verified to share the common
        experiment start (\texttt{2026-08-19T00:00:00Z}) with the GB trace.
\end{enumerate}

The final India trace contains 690 hourly records, covering
\texttt{2026-08-19T00:00:00Z} through \texttt{2026-09-16T17:00:00Z}.
The choice to interpret source timestamps as IST is supported by the physical
timing of solar generation peaks in the source data; however, authoritative
source documentation confirming this convention was not available at the time
of writing. This residual provenance uncertainty is retained as a limitation
in Section~\ref{sec:threats} and does not invalidate the evaluation but calls
for confirmation before final publication. India data also uses average
emissions intensity rather than marginal intensity~\cite{siler2012should};
this distinction is noted but does not change the experimental design.

\subsection{Experiment Start and Temporal Scope}

Both regional traces use a common experiment start of
\texttt{2026-08-19T00:00:00Z}~(UTC). The five-fold time-aware evaluation
partitions the traces into sequential temporal folds following the principles
of blocked cross-validation for time series~\cite{bergmeir2018note}. The folds
are evaluation partitions of a single contiguous trace, not five independent
real-world grid datasets. Reported confidence intervals describe
fold-to-fold variability under this limited design and should not be
interpreted as fully independent replications.

\subsection{Workload Configuration}
\label{subsec:workload}

Experiment~A uses a 60-job workload per run, generated with release slots
sampled from $[0, 7]$ and additional deadline slack from $[2, 13]$ scheduling
slots. Four workload seeds are used per fold, giving 20~paired runs per
region-policy combination across the five folds. Each run uses an identical
workload seed for both regions, enabling direct comparison of scheduling
behaviour under matched job demands against different carbon traces. The
scheduling horizon is 48~slots covering 24~hours.

Experiment~B uses a separate workload generator: release slots are sampled
from $[0, 13]$ and deadline slack from $[2, 20]$ slots. Three seeds are used:
20260320, 20260321, and~20260322. The workload parameters differ intentionally
to explore a wider slack distribution, and the two generators are not directly
comparable. Experiment~B uses India carbon-intensity data starting at
simulation hour~432 of the trace (\texttt{2026-09-06T00:00:00Z}).

\subsection{Scheduling Policies}

Seven scheduling policies are evaluated in Experiment~A.

\textit{FCFS~(baseline).} First-come first-served scheduling allocates jobs in
submission order without reference to carbon intensity. It serves as the
zero-savings reference for all carbon-savings percentages within each
experiment. FCFS admits all 60~submitted jobs; its resource-contention-driven
deadline misses occur because the scheduler makes no attempt to optimise job
ordering with respect to deadlines.

\textit{Bahreini SRR-APX.} This scheduler uses an LP-relaxation formulation
for workload-to-resource allocation, where the objective function incorporates
carbon-intensity-dependent coefficients that favour green-resource utilisation
\cite{bahreini2023dispatcher,bahreini2022approx}. The LP solution is rounded
using K-sample Monte Carlo randomised rounding~\cite{bampis2014energy}, and a
greedy capacity-rejection step removes infeasible allocations. As a consequence
of this rejection step, SRR-APX does not admit all 60~submitted jobs; its
carbon savings and completion-rate figures must be interpreted relative to its
admitted-job count, not the submitted count.

\textit{CADS~(Carbon-Adaptive Deadline-Sensitive), $\theta \in \{0.00, 0.05,
0.10\}$.} CADS parameterises the balance between carbon minimisation and
deadline adherence through threshold~$\theta$. Higher values of~$\theta$ place
more weight on deadline adherence at the cost of smaller carbon savings.

\textit{Deadline-bounded carbon shifting.} This policy shifts workloads to
lower-carbon windows subject to a deadline feasibility constraint applied more
permissively than CADS, prioritising carbon savings at the expense of deadline
compliance.

\textit{Wiesner~2021 Ad-Hoc.} An ad-hoc temporal shifting policy based on the
approach described in~\cite{wiesner2021}, which defers jobs to lower-intensity
windows within their deadline slack. It is included to provide a cross-system
baseline and to characterise its regional sensitivity.

\subsection{Experiment B: Forecast-Guided Variants}
\label{subsec:forecast}

Experiment~B evaluates four variants of the Bahreini SRR-APX scheduler with
and without seasonal-naive forecast guidance~\cite{makridakis2020m4,
hyndman2021forecasting}. The seasonal-naive forecaster repeats the
carbon-intensity observation from the same hour in the preceding period as
its forecast for the current hour.

\begin{itemize}
  \item \textit{Original SRR-APX}: standard Bahreini scheduler with no
        forecast guidance.
  \item \textit{Theta~=~0 verification}: forecast-guided SRR-APX with
        $\theta = 0$, expected to reproduce the original result.
  \item \textit{Forecast-guided $\theta = 0.5$ and $\theta = 1.0$}: the
        forecast carbon signal is blended into the allocation objective via
        Gibbs-style weighting with the specified mixing parameter.
  \item \textit{Oracle-guided diagnostic}: the scheduler is provided with
        future carbon-intensity data. This variant is an offline
        diagnostic reference and is not operationally deployable.
\end{itemize}

This experiment also includes FCFS and Wiesner~2021 Ad-Hoc as Experiment-B
baselines. It does not include CADS or the Deadline-bounded shifting policy.

\subsection{Evaluation Metrics}

The following metrics are reported in both experiments:

\begin{itemize}
  \item \textit{Carbon savings~(\%).} Mean percentage reduction in total
        carbon emissions~(gCO\textsubscript{2}e) relative to the
        experiment-specific FCFS baseline.
  \item \textit{Job admission rate~(\%).} The fraction of the 60~submitted
        jobs that pass the scheduler's capacity-feasibility check and are
        admitted for execution. FCFS and all policies except SRR-APX admit
        all 60~jobs. Admission rate differs from completion rate.
  \item \textit{Completed jobs.} The number of admitted jobs that reach
        execution completion within the simulation horizon.
  \item \textit{Material deadline misses per run.} The mean number of jobs
        whose completion time exceeds the declared deadline by more than
        60~seconds. Jobs with lateness below this threshold are counted as
        any-deadline violations but not material misses.
  \item \textit{Any-deadline violations per run.} All jobs with any positive
        lateness, regardless of magnitude.
  \item \textit{Mean waiting time.} Mean elapsed time from job submission to
        job start across completed jobs. Reported in minutes for
        Experiment~A and seconds for Experiment~B.
  \item \textit{Scheduler runtime~(ms).} Wall-clock duration of the
        scheduling computation per run.
\end{itemize}

\subsection{Audit Validation}

An automated numerical audit suite verified a set of implementation-level
invariants for the simulation pipeline. All tested invariants passed.
Timestamp alignment was not covered by the original audit suite and is
addressed separately through the explicit IST-to-UTC conversion pipeline
described in Section~\ref{subsubsec:india-data}.

%--------------------------------------------------------------
\section{Results}
\label{sec:results}
%--------------------------------------------------------------

\subsection{Cross-Region SRR-APX Comparison}

Table~\ref{tab:crossregion} summarises the Bahreini SRR-APX results for both
regions under the five-fold time-aware evaluation.

\begin{table}[t]
\caption{Cross-Region Bahreini SRR-APX Results\\(Five-Fold, 20~Runs per Region)}
\label{tab:crossregion}
\centering
\begin{tabular}{@{}lcc@{}}
\toprule
\textbf{Metric} & \textbf{Great Britain} & \textbf{India} \\
\midrule
Carbon savings vs.\ FCFS (\%) & 6.04 & 6.16 \\
95\% CI for savings (\%)      & 2.66--9.41 & 4.93--7.38 \\
Material deadline misses/run  & 0.00 & 0.00 \\
Any-deadline violations/run   & 23.05 & 25.20 \\
Job admission rate (\%)       & 94.67 & 93.42 \\
Admitted jobs (of 60)         & 56.80 & 56.05 \\
Completion rate (admitted)    & 100\% & 100\% \\
Mean waiting time (min)       & 129.58 & 163.74 \\
Scheduler runtime (ms)        & 94.71 & 115.98 \\
\bottomrule
\multicolumn{3}{@{}l@{}}{\footnotesize Material miss threshold: lateness $>60$\,s. Completion rate is} \\
\multicolumn{3}{@{}l@{}}{\footnotesize relative to admitted jobs, not to all 60 submitted jobs.} \\
\multicolumn{3}{@{}l@{}}{\footnotesize India timestamp provenance: IST interpretation; see
  Section~\ref{sec:threats}.}
\end{tabular}
\end{table}

The mean carbon savings are 6.04\% in GB and 6.16\% in India, a difference of
0.12~percentage points with India marginally higher. Both regions return zero
material deadline misses on average. SRR-APX admits approximately 94.67\% of
submitted jobs in GB (56.80 of~60) and 93.42\% in India (56.05~of~60); all
admitted jobs complete. The 95\%~confidence intervals overlap substantially,
reflecting fold-to-fold variability within a single contiguous trace rather
than independent replications. The GB interval is notably wider
(2.66\%--9.41\%) than India's (4.93\%--7.38\%), indicating greater
fold-to-fold savings variability in the GB trace. Mean waiting time is
34.16~minutes longer in India, and scheduler runtime is approximately 22\%
longer.

\begin{figure}[t]
    \centering
    \includegraphics[width=\columnwidth]{gb_india_variability_comparison.png}
    \caption{Carbon-intensity variability over the evaluation window for
    Great Britain and India. The comparison reports the standard deviation,
    the 90th--10th percentile range, and the coefficient of variation of
    the hourly carbon-intensity traces used in the scheduling evaluation.
    Table~\ref{tab:crossregion}.}
    \label{fig:grid-variability}
\end{figure}

\subsection{Great Britain: Full Policy Comparison}

Table~\ref{tab:gbpolicies} presents the full GB policy comparison. FCFS has
zero carbon savings by definition and the shortest waiting time~(87.18~min),
but it produces 9.35 material misses per run due to resource contention.
SRR-APX achieves the highest mean savings across all GB policies at 6.04\%,
with zero material misses, at the cost of admitting approximately 94.67\% of
submitted jobs. CADS configurations achieve 4.32\%--5.49\% savings with
4.50--5.30 material misses per run; their CADS~$\theta$=0 variant reaches the
highest savings (5.49\%) among them but also the most material misses (4.85)
and any-violations (29.15). Deadline-bounded shifting achieves 5.50\% savings
but incurs 12.20~material misses. Wiesner~2021 Ad-Hoc achieves only 1.11\%
savings in GB, the lowest among all non-FCFS policies, with 16.15~material
misses per run.

\begin{table*}[t]
\caption{Great Britain: All-Policy Comparison (Five-Fold Evaluation, 20~Runs per Policy).\newline
\footnotesize Savings relative to FCFS. Mat.\,=\,material misses/run (lateness~$>60$\,s).
Any\,=\,any-deadline violations/run. Wait\,=\,mean waiting time. Adm.\,=\,job admission rate.
RT\,=\,scheduler runtime. DB-Shift\,=\,Deadline-bounded carbon shifting.}
\label{tab:gbpolicies}
\centering
\footnotesize
\begin{tabular}{@{}lrrcrrrr@{}}
\toprule
\textbf{Policy} & \textbf{Savings (\%)} & \textbf{95\% CI (\%)} & \textbf{Mat.} & \textbf{Any} & \textbf{Wait (min)} & \textbf{Adm. (\%)} & \textbf{RT (ms)} \\
\midrule
FCFS                  & 0.00 & 0.00--0.00 &  9.35 & 12.30 &  87.18 & 100.0 &  0.00 \\
Bahreini SRR-APX      & 6.04 & 2.66--9.41 &  0.00 & 23.05 & 129.58 &  94.7 & 94.71 \\
CADS $\theta$=0.00    & 5.49 & 1.45--9.53 &  4.85 & 29.15 & 165.81 & 100.0 &  2.01 \\
CADS $\theta$=0.05    & 4.96 & 1.61--8.31 &  4.50 & 25.25 & 149.17 & 100.0 &  1.44 \\
CADS $\theta$=0.10    & 4.32 & 1.20--7.44 &  5.30 & 23.80 & 143.84 & 100.0 &  1.28 \\
DB-Shift              & 5.50 & 0.89--10.10 & 12.20 & 32.70 & 180.17 & 100.0 &  0.90 \\
Wiesner '21 Ad-Hoc    & 1.11 & 0.00--2.22 & 16.15 & 19.50 & 136.94 & 100.0 &  0.96 \\
\bottomrule
\end{tabular}
\end{table*}

\subsection{India: Full Policy Comparison}

Table~\ref{tab:indiapolicies} presents the India policy comparison. SRR-APX
achieves the highest mean savings at 6.16\% with zero material misses,
consistent with its GB performance. However, the service-quality costs of
higher-saving policies are substantially more severe in India than in GB.
Wiesner~2021 Ad-Hoc achieves 5.78\% savings (second-highest in India) but
incurs 35.30~material misses per run and 280.92~minutes of mean waiting,
the worst service-quality outcomes of all policies. Deadline-bounded shifting
achieves 5.27\% savings with 24.80~material misses, also substantially worse
than in GB~(12.20). CADS configurations in India achieve only 4.28\%--4.61\%
savings, noticeably lower than the corresponding GB results, and they produce
15.00--16.30~material misses per run. FCFS again has the shortest waiting time
(87.18~min) with zero savings and 9.35~material misses.

The contrast with GB is notable: Wiesner~2021 Ad-Hoc achieves only 1.11\%
savings in GB but 5.78\% in India, a difference of 4.67~percentage points.
This large region-sensitivity provides direct empirical evidence that cross-region
evaluation is necessary: a policy's savings are not stable across grids.

\begin{table*}[t]
\caption{India: All-Policy Comparison (Five-Fold Evaluation, 20~Runs per Policy).\newline
\footnotesize Savings relative to FCFS. Column definitions match Table~\ref{tab:gbpolicies}.
India IST-to-UTC timestamp alignment described in Section~\ref{subsubsec:india-data}.}
\label{tab:indiapolicies}
\centering
\footnotesize
\begin{tabular}{@{}lrrcrrrr@{}}
\toprule
\textbf{Policy} & \textbf{Savings (\%)} & \textbf{95\% CI (\%)} & \textbf{Mat.} & \textbf{Any} & \textbf{Wait (min)} & \textbf{Adm. (\%)} & \textbf{RT (ms)} \\
\midrule
FCFS                  & 0.00 & 0.00--0.00 &  9.35 & 12.30 &  87.18 & 100.0 &   0.00 \\
Bahreini SRR-APX      & 6.16 & 4.93--7.38 &  0.00 & 25.20 & 163.74 &  93.4 & 115.98 \\
CADS $\theta$=0.00    & 4.61 & 3.55--5.67 & 16.30 & 39.15 & 241.05 & 100.0 &   1.51 \\
CADS $\theta$=0.05    & 4.51 & 3.48--5.53 & 16.00 & 37.40 & 235.83 & 100.0 &   1.10 \\
CADS $\theta$=0.10    & 4.28 & 3.33--5.23 & 15.00 & 32.75 & 219.82 & 100.0 &   1.12 \\
DB-Shift              & 5.27 & 4.28--6.26 & 24.80 & 46.90 & 275.47 & 100.0 &   0.79 \\
Wiesner '21 Ad-Hoc    & 5.78 & 4.83--6.73 & 35.30 & 38.20 & 280.92 & 100.0 &   0.80 \\
\bottomrule
\end{tabular}
\end{table*}

\subsection{Experiment B: Forecast-Guided Investigation}

Table~\ref{tab:forecast} summarises Experiment~B. This experiment uses a
different workload generator, a different FCFS baseline, and three seeds; its
results must not be interpreted as a continuation of Experiment~A.

\begin{table*}[t]
\caption{Experiment~B: Forecast-Guided SRR-APX Results (India, 3~Seeds, Separate Configuration).\newline
\footnotesize Savings relative to this experiment's own FCFS baseline. Wait: seconds (not minutes).
Viol.: deadline violations. Compl.: mean completed jobs (of 60~submitted). RT: scheduler runtime (ms).
Oracle uses future carbon data and is not deployable. Theta~=~0 verification confirms reproduction
of the original Bahreini baseline.}
\label{tab:forecast}
\centering
\footnotesize
\begin{tabular}{@{}lrrrrrrrr@{}}
\toprule
\textbf{Variant} &
  \textbf{Emis. (gCO\textsubscript{2}e)} & \textbf{SD} &
  \textbf{Sav. (\%)} & \textbf{SD} &
  \textbf{Wait (s)} & \textbf{Viol.} & \textbf{Compl.} & \textbf{RT (ms)} \\
\midrule
FCFS                 & 13{,}981.997 & 545.074 & 0.000 & 0.000 & 2{,}581.09 &  5.0 & 60.00 &   0.00 \\
Wiesner '21 Ad-Hoc   & 13{,}464.228 & 667.806 & 3.724 & 1.758 & 13{,}172.96 & 19.0 & 60.00 &   0.00 \\
Original SRR-APX     & 12{,}276.352 & 425.223 & 12.187 & 0.689 & 8{,}822.39 &  7.0 & 56.67 & 205.52 \\
$\theta$=0 verif.    & 12{,}276.352 & 425.223 & 12.187 & 0.689 & 8{,}822.39 &  7.0 & 56.67 & 161.50 \\
FG $\theta$=0.5      & 12{,}216.550 & 422.958 & 12.609 & 1.395 & 8{,}790.59 &  6.0 & 56.33 & 151.58 \\
FG $\theta$=1.0      & 12{,}216.550 & 422.958 & 12.609 & 1.395 & 8{,}790.59 &  6.0 & 56.33 & 168.36 \\
Oracle $\theta$=1.0  & 12{,}216.550 & 422.958 & 12.609 & 1.395 & 8{,}790.59 &  6.0 & 56.33 & 154.96 \\
\bottomrule
\end{tabular}
\end{table*}

The original SRR-APX achieves 12.19\% savings with mean emissions of
12,276.35~gCO\textsubscript{2}e. The forecast-guided $\theta=0.5$ and
$\theta=1.0$ variants both produce 12,216.55~gCO\textsubscript{2}e at 12.61\%
savings, a descriptive improvement of 0.42~percentage points and
approximately 59.80~gCO\textsubscript{2}e over the original baseline. The
theta~=~0 verification confirms that the forecast-guided procedure reproduces
the original result when the forecast weight is zero. Notably, the
oracle-guided diagnostic produces identical aggregate emissions and savings to
the forecast-guided variants~(12,216.55~gCO\textsubscript{2}e, 12.61\%). The
two forecast theta values do not produce distinct aggregate outcomes at the
reported precision.

The 12.61\% savings in Experiment~B must not be compared directly to the
6.04\%/6.16\% figures from Experiment~A. The two experiments use different
workload distributions, deadline slack ranges, seed counts, and evaluation
protocols.

%--------------------------------------------------------------
\section{Discussion}
\label{sec:discussion}
%--------------------------------------------------------------

\subsection{Cross-Region Transferability}

The 0.12~percentage-point difference in mean SRR-APX savings between GB
(6.04\%) and India (6.16\%) is small, with India marginally outperforming GB.
One possible explanation is that SRR-APX responds to relative variations in
carbon intensity within the scheduling window rather than to absolute intensity
levels, so a coal-heavy grid with moderate relative diurnal variation can yield
savings comparable to a more renewable-heavy grid with similar relative
fluctuation. However, this explanation is speculative without direct analysis
of the carbon-intensity time series in each fold. The evidence supports only
the narrower claim that, under the evaluated workload, traces, and experimental
configuration, the measured mean savings were similar across the two regions.

The significantly different confidence interval widths across regions are
informative. GB exhibits a much wider fold-to-fold spread (2.66\%--9.41\%) than
India (4.93\%--7.38\%). This suggests that GB carbon-intensity patterns have
greater intra-week variability across the trace, causing SRR-APX savings to
fluctuate more depending on which portion of the trace falls in each fold.

The most striking cross-regional contrast is the Wiesner~2021 Ad-Hoc policy:
1.11\% savings in GB versus 5.78\% in India, a 4.67~percentage-point gap.
This policy clearly exploits structural patterns that are stronger in the
India trace than in the GB trace. Its effectiveness is therefore not portable
across regions. This finding directly supports the paper's central claim that
cross-region evaluation is necessary.

\subsection{Service-Quality Trade-offs}

The India policy comparison reveals a pronounced and consistent
carbon-savings--service-quality trade-off. The policies with the highest
savings also incur the most material deadline misses and the longest waiting
times.

\begin{figure*}[t]
    \centering
    \includegraphics[width=0.8\textwidth]{algorithm_tradeoff_comparison.pdf}
    \caption{Carbon savings versus mean material deadline misses per run
    for the seven scheduling policies in Great Britain and India.
    Each panel uses the five-fold evaluation results. Lower deadline-miss
    counts and higher carbon savings indicate a more favorable trade-off.}
    \label{fig:tradeoff}
\end{figure*}

Wiesner~2021 Ad-Hoc achieves 5.78\% savings in India with 35.30~material misses
per run. This means that on a typical 60-job run, more than half of all jobs
experience a material deadline violation. An operator running under any
reasonable service-level objective~(SLO) would find this unacceptable,
regardless of the carbon benefit. Deadline-bounded shifting incurs 24.80~misses
per run at 5.27\% savings: nearly 41\% of jobs miss deadlines materially.

CADS configurations in India achieve lower savings~(4.28\%--4.61\%) and still
produce 15.00--16.30~material misses per run, suggesting that their
parameterisation is insufficiently aggressive at preventing misses under India's
carbon trace. Increasing~$\theta$ from 0.00 to 0.10 reduces waiting time from
241.05 to 219.82~minutes but leaves material misses largely unchanged
(16.30~down to 15.00).

Bahreini SRR-APX is the only policy other than FCFS with zero material misses
in India, and it achieves the highest savings of all seven policies at 6.16\%.
This combination of best-in-class savings and zero material misses is the key
service-quality finding of Experiment~A for India. FCFS has the shortest
waiting time~(87.18~min) but delivers no carbon savings and produces 9.35
material misses per run from resource-contention effects alone.

In GB the pattern is similar but less extreme. SRR-APX again achieves the
highest savings~(6.04\%) and zero material misses. All other policies incur
at least 4.50~material misses. Notably, the CADS range in GB (4.32\%--5.49\%)
is higher than in India (4.28\%--4.61\%), suggesting that the CADS
parameterisation is more effective at extracting savings from the GB trace.

\subsection{Admission and Completion Implications}

SRR-APX rejects approximately 3--4 jobs per run due to its greedy
capacity-rejection step. The 6.04\% and 6.16\% savings figures therefore apply
to a slightly smaller admitted workload than the 60~submitted jobs seen by all
other policies. This admission overhead is a genuine trade-off: the rejection
mechanism is part of why SRR-APX achieves zero material misses. Operators should
account for the approximately 6\% job rejection rate when dimensioning workload
pipelines for SRR-APX deployment.

\subsection{Forecast-Guidance Findings}

The forecast-guided variants in Experiment~B yield a descriptive improvement
of 0.42~percentage points over the original SRR-APX baseline. This
improvement translates to approximately 59.80~gCO\textsubscript{2}e per run.
Whether this constitutes a practically significant reduction depends on the
deployment scale. However, several aspects of the result warrant caution.

First, only three seeds were used. The standard deviations for savings
(0.689~pp for original SRR-APX, 1.395~pp for forecast-guided variants) are
large relative to the 0.42~pp descriptive improvement, and no statistical test
is applied. The improvement is a descriptive observation from a small-sample
experiment, not an established finding.

Second, the oracle-guided diagnostic produces identical aggregate outcomes to
both forecast-guided theta values at the reported precision. The oracle uses
future carbon-intensity information, so its inability to outperform the
seasonal-naive forecast variants in aggregate emissions suggests either that
the seasonal-naive forecaster captures most of the exploitable structure in the
India trace for this workload, or that the improvement is bounded by the
workload and resource configuration rather than by forecast accuracy. Neither
explanation is directly established by the available measurements; both are
plausible.

Third, the theta~=~0.5 and theta~=~1.0 variants produce identical aggregate
emissions at the reported precision. This suggests the forecast weighting
saturates before~$\theta = 0.5$, or that the two theta values produce the same
scheduling decisions under the evaluated conditions.

In summary, Experiment~B yields a small descriptive improvement that merits
further investigation with more seeds and rigorous statistical evaluation, but
does not establish robust forecast-guidance benefit.

%--------------------------------------------------------------
\section{Threats to Validity and Limitations}
\label{sec:threats}
%--------------------------------------------------------------

\textit{India timestamp provenance.} The preprocessing pipeline interprets
source CEA timestamps as IST~(Asia/Kolkata) and converts them to UTC. Physical
evidence from solar generation peaks in the source data is consistent with the
IST interpretation, but authoritative source documentation confirming the
timestamp convention was not available at the time of writing. If the source
timestamps represent a different convention, the UTC alignment of the India
trace would be incorrect and would affect all India results. This limitation
calls for explicit confirmation from CEA data documentation before claims about
the India trace are treated as publication-ready.

\textit{Limited trace duration and fold count.} The five-fold evaluation
partitions a single contiguous trace~(approximately 29 days). The folds are
not independent real-world grid datasets. Confidence intervals reflect
fold-to-fold variability within one trace and should not be extrapolated to
general regional behaviour. Longer traces spanning full seasonal cycles would
provide more reliable characterisation.

\textit{Simulated environment.} The evaluation uses a Java-based scheduling
simulation rather than a production cloud. Simulated experiments cannot capture
hardware heterogeneity, network latency, workload inter-arrival variability, or
the operational overhead of a real scheduler. Reported runtimes and waiting
times reflect simulator behaviour.

\textit{Average versus marginal emissions.} Both carbon traces use
average-intensity estimates rather than marginal emissions factors
\cite{siler2012should}. Average intensity can diverge significantly from
marginal intensity, particularly in grids with high renewable penetration. The
savings estimates in this paper are relative to a FCFS baseline computed on
the same intensity measure, limiting but not eliminating the impact of this
distinction.

\textit{Limited policy coverage.} The evaluation covers seven policies. Other
carbon-aware scheduling approaches~\cite{hanafy2024gaia,hanafy2024carbonscaler,
radovanovic2023} are not included. Results reflect behaviour under the tested
policy set and workload configuration.

\textit{Workload generalization.} All Experiment~A runs use 60~jobs with a
fixed release-slot and deadline-slack distribution. Different mixes, job sizes,
or deadline structures may produce different trade-off surfaces. The two
workload generators in this study differ in deadline slack, and their results
are not directly comparable.

\textit{Carbon data imputation.} The India trace uses linear interpolation for
missing hourly observations. Imputed values introduce uncertainty about
whether the interpolated intensities accurately represent grid behaviour during
those periods.

\textit{Forecast experiment sample size.} Experiment~B uses three seeds only.
The standard deviations are large relative to the observed differences, and no
null hypothesis test is conducted. The descriptive improvement of
0.42~percentage points should be treated as exploratory.

\textit{Statistical inference.} The 95\%~confidence intervals reported
for Experiment~A are descriptive fold-level estimates. No formal null hypothesis
test comparing the two regions is conducted. The evidence supports descriptive
observations but not strong inferential claims about regional differences.

%--------------------------------------------------------------
\section{Conclusion and Future Work}
\label{sec:conclusion}
%--------------------------------------------------------------

This paper presented two experimental investigations of carbon-aware cloud
scheduling. Experiment~A provided a controlled, time-aware comparative
evaluation of seven scheduling policies across Great Britain and India using a
five-fold cross-validation scheme, with explicit IST-to-UTC alignment of the
India carbon-intensity trace. Experiment~B explored the effect of seasonal-naive
forecast guidance on the Bahreini SRR-APX scheduler using India data and three
workload seeds.

The central findings are:

\begin{itemize}
  \item For Bahreini SRR-APX, mean carbon savings were 6.04\% in GB and 6.16\%
        in India, a difference of 0.12~percentage points with India marginally
        higher. This similarity does not establish that the scheduler generalises
        equivalently across all grids or workloads.
  \item SRR-APX achieves zero material deadline misses in both regions while
        admitting approximately 93--95\% of submitted jobs. It is the only
        non-FCFS policy with this property.
  \item Cross-regional policy sensitivity is significant: Wiesner~2021 Ad-Hoc
        achieves 1.11\% savings in GB but 5.78\% in India, a 4.67~pp gap.
        This finding directly confirms that policy effectiveness is not portable
        across grids and that cross-region evaluation is necessary.
  \item The India policy comparison shows that higher-saving policies impose
        severe service-quality costs. Wiesner~2021 Ad-Hoc averages 35.30~material
        misses per run at 5.78\% savings; Deadline-bounded shifting averages
        24.80 at 5.27\%. A carbon-saving ranking alone is insufficient to
        assess scheduling quality.
  \item Experiment~B yielded a small descriptive improvement of 0.42~pp from
        forecast guidance over the original SRR-APX baseline. The oracle-guided
        diagnostic produced identical aggregate outcomes to both forecast
        theta values, and the three-seed design precludes strong conclusions.
\end{itemize}

These findings reinforce the principle that carbon-aware schedulers must be
evaluated jointly across carbon reduction and service quality, and that regional
performance must be measured rather than assumed from results in other grids.

Future work should: (1)~confirm India CEA timestamp provenance from authoritative
source documentation and rerun if the IST interpretation proves incorrect;
(2)~evaluate over longer traces spanning multiple seasons;
(3)~extend the comparison to additional grids, particularly in South and Southeast
Asia, sub-Saharan Africa, and Latin America;
(4)~conduct a rigorous statistical comparison of the forecast-guided variants
using a larger seed set;
(5)~investigate the root cause of the saturating oracle behaviour in Experiment~B;
and (6)~characterise the admission-rejection behaviour of SRR-APX across
workload distributions with different deadline tightness.

\bibliographystyle{IEEEtran}
\bibliography{refs}

\end{document}